\documentclass[11pt]{amsart}
\usepackage{amssymb,amsmath,amsthm}
\usepackage[margin=1.1in]{geometry}
\usepackage{hyperref}

\newtheorem{theorem}{Theorem}
\newtheorem{lemma}[theorem]{Lemma}
\newtheorem{proposition}[theorem]{Proposition}
\newtheorem{corollary}[theorem]{Corollary}
\theoremstyle{definition}
\newtheorem{definition}[theorem]{Definition}
\theoremstyle{remark}
\newtheorem{remark}[theorem]{Remark}

\newcommand{\N}{\mathbb{N}}

\title[Lehmer's totient problem: analytic constraints]{Lehmer's totient problem:\\
analytic constraints on a composite solution\\
\normalsize(the problem itself remains open)}
\author{Deep Bhattacharjee}
\date{October 2026}

\begin{document}

\begin{abstract}
Lehmer asked in 1932 whether a composite integer $n$ can satisfy
$\varphi(n)\mid n-1$. This paper does \emph{not} answer that question.
It gives complete, purely analytical proofs, with no computer-assisted
step, of a set of constraints that any composite solution must satisfy.
Let $k$ be the number of prime factors of such an $n$ and
$a=(n-1)/\varphi(n)$. We prove that $n$ is odd, squarefree and
$n\equiv1\pmod{2^{k}}$; that $k\ge10$ when $3\nmid n$; that, when
$3\nmid n$, either some prime factor of $n$ is $\equiv1\pmod 3$ or
$k\ge 223$; that $a\equiv1\pmod 3$ and $k\ge111$ when $3\mid n$; and
that $n<(2k)^{2^{k}}/k$, so that each $k$ admits only finitely many
solutions. Every inequality is reduced to a comparison of explicit
rationals that can be checked by hand. We also show that the product
method behind the lower bounds cannot by itself exclude $k=11$, and we
state precisely where the analytic argument stops. Most of these
statements are weaker than results already in the literature, which we
cite; the contribution is a self-contained analytic account, not a
resolution of the conjecture.
\end{abstract}

\maketitle

\section{Introduction}

Every prime $p$ satisfies $\varphi(p)=p-1$. Lehmer \cite{Leh32} asked
whether any composite $n$ satisfies
\begin{equation}\label{eq:L}
\varphi(n)\mid n-1 ,
\end{equation}
and conjectured that none does. We call a composite solution of
\eqref{eq:L} a \emph{Lehmer number}. Throughout, $n$ denotes a Lehmer
number,
\[
n=p_1p_2\cdots p_k,\qquad p_1<p_2<\cdots<p_k ,
\qquad a=\frac{n-1}{\varphi(n)}\in\N ,
\]
where the factorisation into \emph{distinct} primes is justified by
Lemma~\ref{lem:basic}.

\subsection*{Status of the problem} Lehmer's conjecture is open. Lehmer
showed that a Lehmer number is odd, squarefree and has at least seven
prime factors \cite{Leh32}; Cohen and Hagis raised this to fourteen
\cite{CH80}; Pomerance proved $n<k^{2^{k}}$ \cite{Pom77}, later sharpened
by Burek and \.{Z}mija \cite{BZ19}; and Hagis showed that a Lehmer number
divisible by $3$ has a very large number of prime factors \cite{Hag88}.
The computational work of Mandal, Bhattacharjee and Bhattacharya
\cite{MBB} reports $k\ge16$. None of these, and nothing in this paper,
excludes Lehmer numbers altogether.

\subsection*{What this paper proves} All proofs below are analytic:
every step is either an algebraic identity, an elementary inequality, or
a comparison between two explicit rational numbers that the reader can
verify with pencil and paper. Our results are:

\begin{theorem}\label{thm:main}
Let $n$ be a Lehmer number with $k$ prime factors and
$a=(n-1)/\varphi(n)$.
\begin{enumerate}
\item[(i)] $n$ is odd and squarefree, $\gcd(n,\varphi(n))=1$, $a\ge2$,
and $n\equiv1\pmod{2^{k}}$.
\item[(ii)] If $3\nmid n$, then $k\ge10$.
\item[(iii)] If $3\nmid n$ and no prime factor of $n$ is $\equiv1\pmod
3$, then $a\ge3$ and $k\ge223$. Consequently, if $3\nmid n$ and $k\le222$,
then $n\equiv1\pmod 3$.
\item[(iv)] If $3\mid n$, then $a\equiv1\pmod 3$, $a\ge4$, and $k\ge111$.
\item[(v)] $p_i\le (k-i+1)\,(p_1\cdots p_{i-1}+1)$ for $1\le i\le k$; in
particular $p_1\le 2k$, and
\[
n<\frac{(2k)^{2^{k}}}{k}.
\]
Hence for each $k$ there are at most finitely many Lehmer numbers with
$k$ prime factors.
\end{enumerate}
\end{theorem}

Parts (i)--(iv) are weaker than known bounds; part (v) is a slightly
weaker form of Pomerance's bound with a short self-contained proof. We
include them because together they form a complete analytic treatment
whose every step can be audited without a computer.

\subsection*{What this paper does not prove}
It does not prove Lehmer's conjecture. Section~\ref{sec:gap} explains
the obstruction: the lower bounds above grow at most linearly in the
information used, the upper bound in (v) is doubly exponential in $k$,
and no argument here forces the two to meet. In particular
Proposition~\ref{prop:limit} shows that the product inequality
\eqref{eq:prod}, which drives (ii), is satisfied by an admissible set of
eleven primes, so that method alone cannot reach $k\ge12$.

\section{Basic structure}

\begin{lemma}\label{lem:basic}
Let $n>1$ be composite with $\varphi(n)\mid n-1$. Then
\begin{enumerate}
\item[(a)] $n$ is odd and squarefree;
\item[(b)] $\gcd(n,\varphi(n))=1$; equivalently, $p\nmid q-1$ for all
primes $p,q$ dividing $n$;
\item[(c)] $a=(n-1)/\varphi(n)\ge2$;
\item[(d)] $2^{k}\mid\varphi(n)$ and $n\equiv1\pmod{2^{k}}$.
\end{enumerate}
\end{lemma}

\begin{proof}
Since $\varphi(n)\mid n-1$ and $\gcd(n,n-1)=1$, every common divisor of
$n$ and $\varphi(n)$ divides $n-1$ and $n$, hence equals $1$. This is
(b) once (a) is known, and it already gives squarefreeness: if $p^2\mid
n$ then $p\mid\varphi(n)$, a contradiction. If $n$ were even, then,
being composite and squarefree, $n$ has an odd prime factor $q$, so
$2\mid q-1\mid\varphi(n)$ and $2\mid\gcd(n,\varphi(n))$, a contradiction.
Thus (a). For squarefree $n=p_1\cdots p_k$ we have
$\varphi(n)=\prod(p_i-1)$, so $p_j\mid\varphi(n)$ if and only if $p_j\mid
p_i-1$ for some $i$, giving the second form of (b).

For (c): $a\ge1$, and $a=1$ means $\varphi(n)=n-1$, which forces every
integer in $[1,n-1]$ to be prime to $n$, so $n$ is prime. Hence $a\ge2$.

For (d): each $p_i$ is odd, so $2\mid p_i-1$ and $2^{k}\mid\prod(p_i-1)=
\varphi(n)\mid n-1$.
\end{proof}

The engine of the lower bounds is the following inequality. Since
$n/\varphi(n)=(n-1)/\varphi(n)+1/\varphi(n)$,
\begin{equation}\label{eq:prod}
\prod_{i=1}^{k}\frac{p_i}{p_i-1}=\frac{n}{\varphi(n)}>a\ge2 .
\end{equation}

\begin{definition}
A finite set $S$ of odd primes is \emph{admissible} if $p\nmid q-1$ for
all $p,q\in S$. Write $\Pi(S)=\prod_{p\in S}p/(p-1)$.
\end{definition}

By Lemma~\ref{lem:basic}(b), the prime factors of a Lehmer number form
an admissible set $S$ with $\Pi(S)>a\ge2$.

\begin{lemma}[Monotonicity]\label{lem:mono}
Let $L=\{\ell_1<\ell_2<\cdots\}$ be a set of primes and let $T\subseteq
L$ with $|T|=m$. Then $\Pi(T)\le\prod_{j=1}^{m}\ell_j/(\ell_j-1)$, and for
$m'\le m$ the latter product over the first $m'$ elements is at most the
product over the first $m$.
\end{lemma}

\begin{proof}
The function $x\mapsto x/(x-1)$ is decreasing on $(1,\infty)$ and every
value exceeds $1$. If $T=\{t_1<\cdots<t_m\}$ then $t_j\ge\ell_j$, so
$t_j/(t_j-1)\le\ell_j/(\ell_j-1)$ for each $j$; multiply. The second
claim holds because each omitted factor exceeds $1$.
\end{proof}

\section{Lehmer numbers prime to $3$}

\begin{proposition}\label{prop:k10}
If $n$ is a Lehmer number and $3\nmid n$, then $k\ge10$.
\end{proposition}

\begin{proof}
Suppose $k\le9$. All prime factors of $n$ are at least $5$. We split
into three cases, and in each one bound $\Pi(S)$ by
Lemma~\ref{lem:mono}, applied with a list $L$ that contains every
possible remaining prime factor. In each case the bound is below $2$,
contradicting \eqref{eq:prod}.

\emph{Case A: $5\nmid n$.} All prime factors lie in $L_A=\{7,11,13,
\dots\}$, the primes $\ge7$. By Lemma~\ref{lem:mono},
\[
\Pi(S)\le\frac{7}{6}\cdot\frac{11}{10}\cdot\frac{13}{12}\cdot
\frac{17}{16}\cdot\frac{19}{18}\cdot\frac{23}{22}\cdot\frac{29}{28}\cdot
\frac{31}{30}\cdot\frac{37}{36}=\frac{3212440751}{1791590400}<1.7931 .
\]

\emph{Case B: $5\mid n$, $7\nmid n$.} By admissibility no other prime
factor is $\equiv1\pmod5$, which removes $11,31,41$. The other at most
eight prime factors lie in $\{13,17,19,23,29,37,43,47,\dots\}$, so
\[
\Pi(S)\le\frac54\cdot\frac{13}{12}\cdot\frac{17}{16}\cdot\frac{19}{18}
\cdot\frac{23}{22}\cdot\frac{29}{28}\cdot\frac{37}{36}\cdot\frac{43}{42}
\cdot\frac{47}{46}=\frac{45528350335}{25751126016}<1.7681 .
\]

\emph{Case C: $5\mid n$ and $7\mid n$.} Now the other prime factors
avoid both $1\pmod5$ (removing $11,31,41$) and $1\pmod 7$ (removing
$29,43$). The other at most seven prime factors lie in
$\{13,17,19,23,37,47,53,\dots\}$, so
\[
\Pi(S)\le\frac54\cdot\frac76\cdot\frac{13}{12}\cdot\frac{17}{16}\cdot
\frac{19}{18}\cdot\frac{23}{22}\cdot\frac{37}{36}\cdot\frac{47}{46}\cdot
\frac{53}{52}=\frac{1041947935}{525533184}<1.9827 .
\]
In each case $\Pi(S)<2$, which contradicts \eqref{eq:prod}. Hence
$k\ge10$.
\end{proof}

\begin{remark}
Each fraction above is a product of nine explicit rationals; the reader
can verify the stated bound by multiplying numerators and denominators
(or by multiplying the decimal factors, rounded up, in turn). The
numerators and denominators displayed are in lowest terms.
\end{remark}

\begin{lemma}[A comparison bound]\label{lem:sum}
For $N\ge1$,
\[
\prod_{j=1}^{N}\frac{6j-1}{6j-2}\;\le\;\frac54\left(\frac{6N+1}{7}
\right)^{1/6}.
\]
\end{lemma}

\begin{proof}
For $N=1$ both sides equal $5/4$. For $j\ge2$, $\log\frac{6j-1}{6j-2}=
\log\bigl(1+\frac1{6j-2}\bigr)\le\frac{1}{6j-2}=f(j)$ with
$f(x)=1/(6x-2)$. Since $f$ is convex on $x>1/3$, the Hermite--Hadamard
inequality gives $f(j)\le\int_{j-1/2}^{j+1/2}f(x)\,dx$. Summing over
$2\le j\le N$,
\[
\sum_{j=2}^{N}\log\frac{6j-1}{6j-2}\le\int_{3/2}^{N+1/2}\frac{dx}{6x-2}
=\frac16\log\frac{6N+1}{7}.
\]
Add $\log(5/4)$ for $j=1$ and exponentiate.
\end{proof}

\begin{lemma}\label{lem:twomod3}
Let $Q$ be a set of $N$ primes, each $\equiv2\pmod3$. Then
$\Pi(Q)\le\prod_{j=1}^{N}\frac{6j-1}{6j-2}$.
\end{lemma}

\begin{proof}
An odd prime $\equiv2\pmod 3$ is $\equiv5\pmod6$. The $j$-th smallest
element of $Q$ is therefore at least the $j$-th positive integer
$\equiv5\pmod 6$, namely $6j-1$. Apply the argument of
Lemma~\ref{lem:mono}.
\end{proof}

\begin{proposition}\label{prop:223}
Let $n$ be a Lehmer number with $3\nmid n$. If some prime factor of $n$
is $\equiv1\pmod3$, then $n\equiv1\pmod3$. If no prime factor of $n$ is
$\equiv1\pmod 3$, then $a\ge3$ and $k\ge223$.
\end{proposition}

\begin{proof}
If $p\mid n$ with $p\equiv1\pmod3$, then $3\mid p-1\mid\varphi(n)\mid
n-1$.

Otherwise every $p_i\equiv2\pmod 3$, so $p_i-1\equiv1\pmod3$, whence
$\varphi(n)\equiv1$ and $n\equiv(-1)^{k}\pmod 3$. From $n-1=a\varphi(n)$
we get $a\equiv(-1)^{k}-1\in\{0,1\}\pmod3$, so $a\ne2$ and, by
Lemma~\ref{lem:basic}(c), $a\ge3$. By \eqref{eq:prod},
Lemma~\ref{lem:twomod3} and Lemma~\ref{lem:sum},
\[
3<\Pi(S)\le\frac54\left(\frac{6k+1}{7}\right)^{1/6},
\qquad\text{so}\qquad
\frac{6k+1}{7}>\left(\frac{12}{5}\right)^{6}=\frac{2985984}{15625}.
\]
For $k\le222$ the left side is at most $1333/7<190.43$, while the right
side exceeds $191.10$. Hence $k\ge223$.
\end{proof}

Propositions~\ref{prop:k10} and~\ref{prop:223} give parts (ii) and (iii)
of Theorem~\ref{thm:main}.

\section{Lehmer numbers divisible by $3$}

\begin{proposition}\label{prop:three}
If $n$ is a Lehmer number and $3\mid n$, then $a\equiv1\pmod3$, $a\ge4$,
and $k\ge111$.
\end{proposition}

\begin{proof}
By Lemma~\ref{lem:basic}(b), no other prime factor $q$ of $n$ is
$\equiv1\pmod3$, so each such $q$ satisfies $q\equiv2\pmod3$ and
$q-1\equiv1\pmod 3$. Hence $\varphi(n)=2\prod_{q\ne3}(q-1)\equiv2\pmod
3$, while $n-1\equiv-1\equiv2\pmod 3$. From $n-1=a\varphi(n)$ we get
$2a\equiv2$, so $a\equiv1\pmod3$; with $a\ge2$ this forces $a\ge4$.

Let $Q$ be the set of the $k-1$ prime factors other than $3$. By
\eqref{eq:prod}, $\frac32\Pi(Q)>4$, i.e.\ $\Pi(Q)>8/3$. By
Lemmas~\ref{lem:twomod3} and~\ref{lem:sum} with $N=k-1$,
\[
\frac83<\frac54\left(\frac{6N+1}{7}\right)^{1/6},
\qquad\text{so}\qquad
\frac{6N+1}{7}>\left(\frac{32}{15}\right)^{6}=\frac{1073741824}
{11390625}>94.26 .
\]
For $N\le109$ the left side is at most $655/7<93.58$. Hence $N\ge110$
and $k\ge111$.
\end{proof}

\section{An explicit upper bound}

\begin{lemma}\label{lem:prefix}
Let $n=p_1\cdots p_k$ be a Lehmer number and $1\le i\le k$. Put
$m=p_1\cdots p_{i-1}$ (so $m=1$ when $i=1$) and $r=n/m=p_i\cdots p_k$.
Then $m/\varphi(m)<a$, and consequently
\[
\frac{r}{\varphi(r)}>1+\frac1m .
\]
\end{lemma}

\begin{proof}
Suppose $m/\varphi(m)\ge a$. Since $r>1$ and $r$ is divisible by $p_i$,
$r/\varphi(r)\ge p_i/(p_i-1)=1+1/(p_i-1)\ge 1+1/(p_k-1)$. Then
\[
\frac1{\varphi(n)}=\frac{n}{\varphi(n)}-a
=\frac{m}{\varphi(m)}\cdot\frac{r}{\varphi(r)}-a
\ge a\Bigl(\frac{r}{\varphi(r)}-1\Bigr)\ge\frac{2}{p_k-1},
\]
so $2\varphi(n)\le p_k-1$. But $\varphi(n)=(p_k-1)\varphi(n/p_k)\ge
p_k-1$, a contradiction. Hence $m<a\varphi(m)$, and since both are
integers, $a\varphi(m)\ge m+1$. Therefore
\[
\frac{r}{\varphi(r)}=\frac{n/\varphi(n)}{m/\varphi(m)}>\frac{a\varphi(m)}
{m}\ge 1+\frac1m. \qedhere
\]
\end{proof}

\begin{proposition}\label{prop:upper}
With the notation of Lemma~\ref{lem:prefix}, $p_i\le(k-i+1)(m+1)$. In
particular $p_1\le2k$, and $n<(2k)^{2^{k}}/k$.
\end{proposition}

\begin{proof}
Let $s=k-i+1$, the number of primes dividing $r$, each at least $p_i$.
By Lemma~\ref{lem:prefix} and the inequalities $1+x\le e^{x}$ and
$\log(1+1/m)\ge1/(m+1)$,
\[
\exp\Bigl(\frac{1}{m+1}\Bigr)\le1+\frac1m<\frac{r}{\varphi(r)}\le
\Bigl(1+\frac{1}{p_i-1}\Bigr)^{s}\le\exp\Bigl(\frac{s}{p_i-1}\Bigr).
\]
Hence $p_i-1<s(m+1)$, and as $p_i$ is an integer, $p_i\le s(m+1)$. For
$i=1$ this reads $p_1\le2k$.

Put $x_i=p_1\cdots p_i$, so $x_0=1$ and $x_k=n$. The bound just proved
gives $x_i=x_{i-1}p_i\le k\,x_{i-1}(x_{i-1}+1)$, hence
\[
x_i+1\le k\,x_{i-1}(x_{i-1}+1)+1\le k\,(x_{i-1}+1)^{2}.
\]
With $y_i=k(x_i+1)$ this is $y_i\le y_{i-1}^{2}$, so $y_k\le
y_0^{2^{k}}=(2k)^{2^{k}}$, that is $n=x_k<(2k)^{2^{k}}/k$.
\end{proof}

\begin{proof}[Proof of Theorem~\ref{thm:main}]
Part (i) is Lemma~\ref{lem:basic}; (ii) is Proposition~\ref{prop:k10};
(iii) is Proposition~\ref{prop:223}; (iv) is Proposition~\ref{prop:three};
(v) is Proposition~\ref{prop:upper}. Finiteness for each $k$ follows
because $n$ is bounded in terms of $k$.
\end{proof}

\section{Where the analytic method stops}\label{sec:gap}

We record exactly what the arguments above cannot do.

\begin{proposition}\label{prop:limit}
The set $S=\{5,7,13,17,19,23,37,59,67,73,83\}$ is admissible and
\[
\Pi(S)=\frac{2995513256722805}{1484664878923776}>2 .
\]
Consequently the conditions used in Section~3 (admissibility and
inequality \eqref{eq:prod} with $a=2$) do not exclude $k=11$.
\end{proposition}

\begin{proof}
Admissibility: the numbers $p-1$ for $p\in S$ are $4,6,12,16,18,22,36,58,
66,72,82$, whose odd prime factors are among $3,11,29,41$, none of which
lies in $S$, and $2\notin S$. For the product, $2\cdot1484664878923776=
2969329757847552<2995513256722805$.
\end{proof}

Proposition~\ref{prop:limit} does not produce a Lehmer number: the
product $n'=\prod_{p\in S}p$ need not satisfy $\varphi(n')\mid n'-1$,
and in fact the method never uses the divisibility \eqref{eq:L} beyond
the two consequences $\gcd(n,\varphi(n))=1$ and $n/\varphi(n)>2$ (plus
congruences modulo $3$). Any proof of Lehmer's conjecture must use
\eqref{eq:L} in an essentially stronger way. Concretely, the gaps are:

\begin{enumerate}
\item[(G1)] \emph{No evident stopping point.} Proposition~\ref{prop:limit}
shows that admissibility and \eqref{eq:prod} are compatible with $k=11$.
We see no reason why they should become incompatible for larger $k$,
but we do not prove that they remain compatible for every $k$; we
record this only to explain why the method, as it stands, cannot
terminate in a proof of the conjecture.
\item[(G2)] \emph{Upper and lower bounds do not meet.} Theorem
\ref{thm:main}(v) bounds $n$ doubly exponentially in $k$, while the
lower bounds are bounds on $k$ alone. Combining them yields only that
$n$ lies in a finite set for each $k$, not that the set is empty.
\item[(G3)] \emph{Finite but unchecked ranges.} For each fixed $k\ge10$
(with $3\nmid n$) the problem reduces, by (v), to a finite search. That
search is computational and lies outside the scope of this paper, which
admits no computer-assisted step.
\end{enumerate}

\noindent Lehmer's totient problem therefore remains open, and this paper
claims no more than Theorem~\ref{thm:main}.

\begin{thebibliography}{9}

\bibitem{BZ19} D.~Burek and B.~\.{Z}mija, \emph{A new upper bound for
numbers with the Lehmer property and its application to repunit
numbers}, Int.\ J.\ Number Theory \textbf{15} (2019).

\bibitem{CH80} G.~L.~Cohen and P.~Hagis, Jr., \emph{On the number of
prime factors of $n$ if $\varphi(n)\mid n-1$}, Nieuw Arch.\ Wisk.\ (3)
\textbf{28} (1980), 177--185.

\bibitem{Hag88} P.~Hagis, Jr., \emph{On the equation $M\varphi(n)=n-1$},
Nieuw Arch.\ Wisk.\ (4) \textbf{6} (1988), 255--261.

\bibitem{Leh32} D.~H.~Lehmer, \emph{On Euler's totient function}, Bull.\
Amer.\ Math.\ Soc.\ \textbf{38} (1932), 745--751.

\bibitem{MBB} P.~Mandal, D.~Bhattacharjee and U.~Bhattacharya,
\emph{Lehmer's totient problem with fewer than sixteen prime factors},
preprint; code and data at \url{https://github.com/creelie/Lehmer-1-2}.

\bibitem{Pom77} C.~Pomerance, \emph{On composite $n$ for which
$\varphi(n)\mid n-1$, II}, Pacific J.\ Math.\ \textbf{69} (1977),
177--186.

\end{thebibliography}

\end{document}
